\section{Markings, coefficients, and conventions}\label{sec:setup}

The detector in Section~\ref{sec:detector} gives a condition that any finite
construction must satisfy. We now prepare the coefficient calculation that
will violate it. The prime remains \(p=3\), and \(k\) is the finite field
chosen there, containing \(\mathbb F_{p^6}\) and enlarged finitely when
needed to fix the constant identifications. Recall the Honda groups
\(\Gamma_i\), their division algebras \(D_i\), and their stabilizers
\(P_i=\mathcal O_{D_i}^{\times}\). The Honda forms can be chosen over
\(\mathbb F_p\), and their endomorphisms are defined over
\(\mathbb F_{p^i}\). Thus \(k\) contains both endomorphism fields needed
in heights two and three. Set
\begin{equation}\label{eq:stabilizers}
 P=P_3,\qquad
 H=\ker(\Nrd:P\longrightarrow\Zp^\times),\qquad
 J=P_2.
\end{equation}
The actions below are the unextended stabilizer actions: they fix \(k\).
When Galois descent is involved, it will be handled explicitly rather
than included silently in \(P\) or \(J\).

Choose the unramified quadratic subfield
\(F\subset D_2\). The character used in the obstruction is
\begin{equation}\label{eq:norm-character}
 \delta:\mathcal O_F^\times\longrightarrow\mathbb F_3^\times,
 \qquad
 \delta(g)=\overline{N_{F/\mathbb Q_3}(g)}.
\end{equation}
It is nontrivial because the residue-field norm
\(\mathbb F_9^\times\to\mathbb F_3^\times\) is surjective.
Its inverse and its unramified Galois conjugate are equal to it.

We also use
\[
 H'=\Nrd^{-1}(\mu_2),\qquad
 H'/H=\mu_2,\qquad
 P/H'\cong 1+3\mathbb Z_3\cong\mathbb Z_3.
\]
Here the reduced norm on maximal-order units is surjective. For these
particular algebras, one may see this on the units of an unramified
maximal subfield: the residue norm is surjective, and the norm on
principal units is surjective by the trace and the \(3\)-adic
logarithm. No group section of the full reduced norm is needed.

\subsection{The one-parameter stratum}

We first fix the coordinate normalization behind the Hasse coefficients.
Lubin--Tate construct a universal deformation from a normalized
representative of the special fibre \citep[Proposition~1.1, p.~50,
and Theorem~3.1]{LT66}. On the stratum where the preceding parameters
vanish, their parameter \(\tau_r\), with \(q=p^r\), occurs as
\[
 \Phi(X,Y)\equiv X+Y+\tau_r C_q(X,Y)
       \pmod{\text{total degree }q+1},
 \qquad
 C_q(X,Y)=\frac{(X+Y)^q-X^q-Y^q}{p}.
\]
The expression defining \(C_q\) is the integral polynomial used in
\citet[equations~(2.4)--(2.5), p.~256]{Lazard55}. Iterating the group
law \(p\) times and reducing in characteristic \(p\) gives
\[
 \begin{split}
 [p]_{\Phi}(T)
 &\equiv \tau_r\sum_{j=1}^{p-1} C_q(jT,T)\\
 &=\tau_r\frac{p^q-p}{p}T^q
  \equiv-\tau_r T^q\pmod{T^{q+1}}.
 \end{split}
\]
Thus the raw successive Hasse coefficient in that normalized
representative is \(-\tau_r\).

The Honda coordinate was already fixed in Section~\ref{sec:detector}.
Choose the special-fibre coordinate isomorphism from the normalized
representative to that Honda law and transport the universal family
along it. Write \(\lambda_h\in k^\times\) for its linear coefficient
in height \(h\). The normalization preceding Lubin--Tate Proposition~1.1
uses an isomorphism over the original residue field. Since the height-two
Honda law is defined over \(\mathbb F_9\), we perform that construction
already there, so \(\lambda_2\in\mathbb F_9^\times\). This transport makes the special fibre literally the
chosen Honda law, so its \(p\)-series is still \(T^{p^h}\). Under a
coordinate change \(T'=cT+O(T^2)\), the first nonzero coefficient at
\(T^{p^r}\) becomes \(c^{1-p^r}\) times the old coefficient. We
therefore name the height-three parameters
\[
 a=-\lambda_3^{1-p}\tau_1,\qquad
 t=-\lambda_3^{1-p^2}\tau_2,
\]
and name the height-two parameter
\(x=-\lambda_2^{1-p}\tau_1\). These are unit rescalings of regular
parameters. In the transported families the first raw Hasse
coefficient is exactly \(a\), the next one modulo \(a\) is exactly
\(t\), and the first height-two coefficient is exactly \(x\). These
are chosen coordinate normalizations, rather than identities
independent of a coordinate. In the Lubin--Tate representative itself
they read \(a=-\tau_1\), \(t=-\tau_2\), and \(x=-\tau_1\).

We choose the mixed-characteristic height-two parameter compatibly.
Lift the chosen height-two coordinate isomorphism coefficientwise to
\(W(\mathbb F_9)\); its linear coefficient remains a unit. Transport the
mixed-characteristic Lubin--Tate family along this lift. The resulting
family over \(W(\mathbb F_9)[[\tau_1]]\) has a regular
parameter lifting this \(x\): multiply \(\tau_1\) by a lift of the
nonzero constant \(-\lambda_2^{1-p}\). We denote the lift by \(x\) as
well and use the scalar-extended family over \(W(k)[[x]]\) for \(E_2(k)\).
Functorial Morava \(E\)-theory realization
\citep[Section~7 and Corollary~7.7]{GH04} supplies the map from
\(E_2(\mathbb F_9)\) to \(E_2(k)\); its coefficient map sends the chosen
regular parameter \(x\) to \(x\). This convention will be
used in the completion comparison in Section~\ref{sec:homotopy}.

Let
\[
                    A=k[[a,t]]
\]
be the universal equal-characteristic deformation ring of
\(\Gamma_3\) in this normalization \citep[Theorem~3.1]{LT66}.
Let \(\mathcal G\) denote the associated algebraic \(p\)-divisible
group over \(A\).

We specify the order of construction, since it matters after \(t\)
is inverted. The formal multiplication-by-\(p^r\) series over \(A\)
is distinguished of degree \(p^{3r}\).
Weierstrass preparation gives its finite flat kernel over \(A\);
the formal group law supplies its group structure. These kernels,
with their compatible inclusions, form \(\mathcal G\).
This is the finite-torsion construction over a complete local base in
\citet[Section~2.2, Proposition~1]{Tate67}.
This equivalence also transfers universality to the algebraic family.
Let \(B\) be a complete Noetherian local \(k\)-algebra with residue field
\(k\), and let \(H/B\) be a marked deformation of
\(\Gamma_3[p^\infty]\). For \(Q_r=\mathcal O(H[p^r])\), the algebra
\[
 Q_r/\mathfrak m_BQ_r\cong k[T]/(T^{p^{3r}})
\]
is local. Since \(Q_r\) is finite over \(B\), every maximal ideal lies
over \(\mathfrak m_B\), so \(Q_r\) is local and \(H\) is connected in
Tate's sense. The inverse equivalence gives a formal group of dimension
one, as on the closed fibre; a coordinate lifting the marking makes it
a Lubin--Tate deformation. Full faithfulness identifies its marked isomorphisms
\citep[Section~2.2, Proposition~1, and Section~2.3]{Tate67}.
The prepared finite kernels commute with local base change. Hence
Lubin--Tate universality transfers to the above algebraic
\(\mathcal G/A\), supplying the universality used by the later
Kodaira--Spencer argument \citep[Theorem~3.1 and paragraph~3.2]{LT66}.
We then pull this finite torsion system back along
\[
                    A\longrightarrow K=k((t)),\qquad a\longmapsto0.
\]
On this stratum the \(p\)-series factors through the
\(p^2\)-power Frobenius, and the remaining homomorphism has invertible
linear coefficient \(t\). Thus \(\mathcal G_K\) has a connected
part of height two and an \'etale quotient of height one.
Forming only the formal completion at the identity after this base
change would discard the latter quotient.

The ring-theoretic differential statement needed later is elementary
in these coordinates.

\begin{lemma}\label{lem:finite-p-basis}
The ordinary module \(\Omega^1_{A/k}\) is free on \(da,dt\) and agrees
with the module of continuous differentials.
The element \(t\) is a \(p\)-basis of \(K\), remains a \(p\)-basis at
each finite separable field extension of \(K\), and gives
\(\Omega^1_{L/k}=L\,dt\) for the marking algebra defined below.
\end{lemma}

\begin{proof}
Since \(k\) is perfect, every \(f\in A\) has a unique expansion
\[
          f=\sum_{0\leq i,j<p}a^it^j f_{ij}^{\,p},
          \qquad f_{ij}\in A.
\]
Ordinary derivations therefore have the usual partial-derivative
formula and are determined freely by the values of \(a,t\).
The continuous partial derivatives show that these values are
independent. This proves the assertion for \(A\).
The corresponding one-variable expansion proves the assertion for
\(K\). A \(p\)-basis is preserved by a separable algebraic extension:
Frobenius base change along a separable extension is an isomorphism,
so the basis of \(K\) over \(K^p\) remains a basis after that base
change. The differential statement follows by the same argument.
For a finite \'etale algebra it holds on each field factor, and it
passes to the filtered union defining \(L\).
\end{proof}

The universal deformation theory of \(p\)-divisible groups gives a
functorial Kodaira--Spencer isomorphism over \(A\)
\citep[Section~5]{Lau10}; the version over
arbitrary square-zero target rings is recalled where it is used in
Proposition~\ref{prop:invariant-differential} and
Theorem~\ref{thm:jet-comparison}.
Lemma~\ref{lem:finite-p-basis} explains why the ordinary algebraic
differentials in that theorem are available after specializing to
the Laurent field \(K\).

\subsection{The connected marking tower}

The connected-marking tower is the coheight-one construction studied by
Gross and Torii \citep[Theorem~3.5(1a)]{Gross79}
\citep[Section~2.3 and Theorem~2.9]{Torii03}. In Torii's Laurent-field
setting the lower-height stabilizer acts simply transitively on markings
and commutes with the higher-height action. We use the ind-\'etale torsor
form, allowing products of fields at finite stages, and prove the integral
\'etale marking and finite-level descent needed below.

Let \(L\) be the algebra classifying isomorphisms between the
connected formal law of \(\mathcal G_K\) and the constant Honda
formal group \(\Gamma_2\). Its coefficient construction and
separability are proved in Section~\ref{sec:towers}. The marking
space is a pro-\'etale torsor under the pro-constant group \(J\).
We write
\begin{equation}\label{eq:marking-algebra}
          L=\colim_{U\subset J} L^U ,
\end{equation}
where \(U\) runs through open subgroups and \(L^U\) is the finite
\'etale \(K\)-algebra of markings modulo \(U\).
The superscript is compatible with the \(J\)-action on the tower.
The finite stages need not be fields.

The height-three stabilizer \(P\) acts on the universal deformation and
preserves its height strata \citep[Proposition~3.3 and paragraph~3.4]{LT66}.
It also transports a connected marking.
Changes of the connected marking give the commuting \(J\)-action.
Thus \(P\) preserves each \(L^U\), while \(J\) carries stages to
the corresponding conjugate stages.
The actions on finite jets and their \'etale lifts are continuous.
Normalize the valuation of \(K\) by \(v(t)=1\), extend it to each
field factor of \(L^U\), and give a product of factors its product
topology. The \(P\)-action preserves these valuations, since its
action on the height-two stratum takes \(t\) to a unit multiple of
\(t\).

For a finite stage \(B=L^U\), put
\[
 B_e=B^{1/p^e},\qquad
 B^{\mathrm{perf}}=\bigcup_{e\geq0}B_e,\qquad
 B^\flat=\widehat{B^{\mathrm{perf}}}.
\]
Roots are taken inside the perfection and all valuations are
extended from \(B\). Each \(B_e\) is complete. The symbol \(B^\flat\)
in this paper means the \emph{completed perfection of this finite
stage}; it does not denote the completion of the entire marking
tower. The other perfection used throughout the proof is
\[
                         R=L^{\mathrm{perf}}
                          =\colim_{U,e}(L^U)^{1/p^e}.
\]
These conventions are applied factorwise to products of fields.

\subsection{Continuous cochains at finite stages}

We use inhomogeneous continuous group cochains. If a compact group
\(Q\) acts continuously on a complete coefficient module \(M\), then
\(C^s_{\mathrm{cts}}(Q,M)=
\operatorname{Map}_{\mathrm{cts}}(Q^s,M)\), with the usual bar
differential. An extra compact profinite parameter space \(T\)
means
\[
 C^s_{\mathrm{cts}}(Q,M;T)
       =\operatorname{Map}_{\mathrm{cts}}(T\times Q^s,M);
\]
the differential leaves \(T\) fixed.
The parameter space need not be metrizable.

\begin{definition}[Stage cochains]\label{def:stage-cochains}
For \(Q=P,H,H'\), or a closed subgroup of one of these groups,
cochains with the incomplete coefficients \(L\) and \(R\) mean
\begin{align}
 C^s_{\mathrm{cts}}(Q,L;T)
   &:=\colim_U
       \operatorname{Map}_{\mathrm{cts}}(T\times Q^s,L^U),
       \nonumber\\
 C^s_{\mathrm{cts}}(Q,R;T)
   &:=\colim_{U,e}
       \operatorname{Map}_{\mathrm{cts}}
            (T\times Q^s,(L^U)^{1/p^e}).
       \label{eq:stage-cochains}
\end{align}
We also use the completed-stage complex
\begin{equation}\label{eq:completed-stage-cochains}
 \colim_U C^\bullet_{\mathrm{cts}}(Q,(L^U)^\flat;T).
\end{equation}
All filtered colimits in these definitions are taken degreewise.
Omitting \(T\) means that \(T\) is a point.
\end{definition}

Thus each element of one of the complexes in
Definition~\ref{def:stage-cochains} has its values in a single complete
finite marking or root stage. We do not identify this convention
with all continuous maps into an incompletely topologized union.
The algebraic inclusion \(L\to R\), and the inclusions of finite
root stages into their completed perfections, give natural maps
between these complexes. They retain the \(P\)- and \(J\)-actions
on the full directed systems. A proof may restrict to cofinally
many stages containing specified finite data; the maps themselves
are defined before that restriction.

The use of parameters is part of each comparison theorem below.
It supplies the iterated cochains for quotient descent and the
continuous function coefficients of Morava cooperations. On a
discrete target, compactness makes continuous functions locally
constant with finite image. On a complete valuation target it
instead gives uniform approximation on finite clopen partitions;
the distinction will be used in
Section~\ref{sec:decompletion}.

\subsection{Geometric notation}

For a perfectoid base in characteristic \(p\), let \(V_i\) denote
the basic vector bundle of rank \(i\) and degree one on the
relative Fargues--Fontaine curve. The notation
\[
                  N=H^1(V_2^\vee),\qquad N^*=N\setminus\{0\}
\]
means the associated cohomology sheaf and its complement of the
zero section on the relevant \(v\)-site.
Relative section-sheaf statements, rather than just statements
about individual geometric fibers, will be used in
Section~\ref{sec:towers}. Over an algebraically closed perfectoid
field \(C^\flat\), a choice of untilt \(C\) gives the concrete
presentation of \(N^*\) used for the analytic calculation.
Base Frobenius is defined before this choice.

The height-two Morava theory and its periodicity line used by the homotopy
test were fixed in Section~\ref{sec:detector}. The height-three theory
will enter after the ordinary coefficient calculation is complete.
